% GENERATED FILE. Edit canonical lessons, then run npm run generate:books.
% canonical-source-hash: fnv1a64:28fd3e8dcc093d1a
% canonical-lessons: AR-C419-ijabi, AR-C419-salbi, AR-C419-asasi, AR-C419-thanawi, AR-C419-shamil

\chapter{Positive and negative}
\label{ch:ar-positive-and-negative}
\hlchaptermodality{419}

\begin{chapteropening}
\textbf{By the end of this chapter:} \emph{I can say positive, negative, basic, secondary and comprehensive.}

Complete the last lesson of chapter 419: I can say positive, negative, basic, secondary and comprehensive.
\end{chapteropening}

\section[\texorpdfstring{\ar{إيجابي} (ījābī)}{إيجابي (ījābī)}]{\ar{إيجابي} (ījābī) — positive}

\label{lesson:AR-C419-ijabi}



\begin{quote}
Before the new one: say the Arabic for effectiveness, then the Arabic for independent.
\end{quote}

\subsection*{You'll want to know: \ar{إيجابي}}
\textbf{\ar{إيجابي}} — \emph{ījābī} — \textquotedblleft{}positive\textquotedblright{}.

\begin{quote}
\ar{كان} \ar{للقرار} \ar{أثر} \ar{إيجابي}.
\end{quote}

\begin{quote}
\emph{kāna li-l-qarār athar ījābī.}
\end{quote}

\begin{quote}
The decision had a positive effect.
\end{quote}

\begin{cousinweb}
\textbf{\ar{إيجابي}} (\emph{ījābī}) is built on the native Arabic root \emph{w-j-b}, \textquotedblleft{}to be due, binding\textquotedblright{} (Lane). \emph{Ījābī} is a modern adjective from \emph{ījāb}, the Form IV verbal noun, affirming: positive.
\end{cousinweb}

\begin{grammarlens}[title={What you've built}]
One more word for saying what happens if.
\end{grammarlens}

\subsection*{Your turn}
\begin{itemize}
\raggedright
  \item \emph{Say it:} \emph{ījābī}
  \item \emph{Say it:} \emph{ījābī}, once more
  \item \emph{Say it:} \emph{kāna li-l-qarār athar ījābī.} — The decision had a positive effect.
  \item \emph{Recall:} say the Arabic for effectiveness, then the Arabic for independent, then say \emph{ījābī} again
\end{itemize}

\begin{tcolorbox}[breakable,colback=teal!4,colframe=teal!35!black,title={Before you move on}]
What does \ar{إيجابي} mean? (\textquotedblleft{}Positive\textquotedblright{}.) Say it once more.
\end{tcolorbox}
\section[\texorpdfstring{\ar{سلبي} (salbī)}{سلبي (salbī)}]{\ar{سلبي} (salbī) — negative}

\label{lesson:AR-C419-salbi}



\begin{quote}
Before the new one: say the Arabic for independent, then the Arabic for positive.
\end{quote}

\subsection*{You'll want to know: \ar{سلبي}}
\textbf{\ar{سلبي}} — \emph{salbī} — \textquotedblleft{}negative\textquotedblright{}.

\begin{quote}
\ar{لا} \ar{يحب} \ar{أحد} \ar{الشخص} \ar{السلبي}.
\end{quote}

\begin{quote}
\emph{lā yuḥibbu aḥad ash-shakhṣ as-salbī.}
\end{quote}

\begin{quote}
Nobody likes a negative person.
\end{quote}

\begin{cousinweb}
\textbf{\ar{سلبي}} (\emph{salbī}) is built on the native Arabic root \emph{s-l-b}, \textquotedblleft{}to strip, to take away\textquotedblright{} (Lane). \emph{Salbī} is a modern adjective from \emph{salb}, taking away: negative.
\end{cousinweb}

\begin{grammarlens}[title={What you've built}]
One more word for saying what happens if.
\end{grammarlens}

\subsection*{Your turn}
\begin{itemize}
\raggedright
  \item \emph{Say it:} \emph{salbī}
  \item \emph{Say it:} \emph{salbī}, once more
  \item \emph{Say it:} \emph{lā yuḥibbu aḥad ash-shakhṣ as-salbī.} — Nobody likes a negative person.
  \item \emph{Recall:} say the Arabic for independent, then the Arabic for positive, then say \emph{salbī} again
\end{itemize}

\begin{tcolorbox}[breakable,colback=teal!4,colframe=teal!35!black,title={Before you move on}]
What does \ar{سلبي} mean? (\textquotedblleft{}Negative\textquotedblright{}.) Say it once more.
\end{tcolorbox}
\section[\texorpdfstring{\ar{أساسي} (asāsī)}{أساسي (asāsī)}]{\ar{أساسي} (asāsī) — basic}

\label{lesson:AR-C419-asasi}



\begin{quote}
Before the new one: say the Arabic for positive, then the Arabic for negative.
\end{quote}

\subsection*{You'll want to know: \ar{أساسي}}
\textbf{\ar{أساسي}} — \emph{asāsī} — \textquotedblleft{}basic\textquotedblright{}.

\begin{quote}
\ar{الخبز} \ar{طعام} \ar{أساسي} \ar{هنا}.
\end{quote}

\begin{quote}
\emph{al-khubz ṭaʿām asāsī hunā.}
\end{quote}

\begin{quote}
Bread is a basic food here.
\end{quote}

\begin{cousinweb}
\textbf{\ar{أساسي}} (\emph{asāsī}) is built on the native Arabic root \emph{ʾ-s-s}, \textquotedblleft{}to lay a foundation\textquotedblright{} (Lane). \emph{Asāsī} is the adjective of \emph{asās}, a foundation: basic, fundamental.
\end{cousinweb}

\begin{grammarlens}[title={What you've built}]
One more word for saying what happens if.
\end{grammarlens}

\subsection*{Your turn}
\begin{itemize}
\raggedright
  \item \emph{Say it:} \emph{asāsī}
  \item \emph{Say it:} \emph{asāsī}, once more
  \item \emph{Say it:} \emph{al-khubz ṭaʿām asāsī hunā.} — Bread is a basic food here.
  \item \emph{Recall:} say the Arabic for positive, then the Arabic for negative, then say \emph{asāsī} again
\end{itemize}

\begin{tcolorbox}[breakable,colback=teal!4,colframe=teal!35!black,title={Before you move on}]
What does \ar{أساسي} mean? (\textquotedblleft{}Basic\textquotedblright{}.) Say it once more.
\end{tcolorbox}
\section[\texorpdfstring{\ar{ثانوي} (thānawī)}{ثانوي (thānawī)}]{\ar{ثانوي} (thānawī) — secondary}

\label{lesson:AR-C419-thanawi}



\begin{quote}
Before the new one: say the Arabic for negative, then the Arabic for basic.
\end{quote}

\subsection*{You'll want to know: \ar{ثانوي}}
\textbf{\ar{ثانوي}} — \emph{thānawī} — \textquotedblleft{}secondary\textquotedblright{}.

\begin{quote}
\ar{هذه} \ar{مسألة} \ar{ثانوية}.
\end{quote}

\begin{quote}
\emph{hādhihi masʾala thānawiyya.}
\end{quote}

\begin{quote}
This is a secondary matter.
\end{quote}

\ar{مدرسة} \ar{ثانوية}, a secondary school.

\begin{cousinweb}
\textbf{\ar{ثانوي}} (\emph{thānawī}) is built on the native Arabic root \emph{th-n-y}, \textquotedblleft{}to fold; to make two\textquotedblright{} (Lane). \emph{Thānawī} is a modern adjective from \emph{thānī}, second: secondary.
\end{cousinweb}

\begin{grammarlens}[title={What you've built}]
One more word for saying what happens if.
\end{grammarlens}

\subsection*{Your turn}
\begin{itemize}
\raggedright
  \item \emph{Say it:} \emph{thānawī}
  \item \emph{Say it:} \emph{thānawī}, once more
  \item \emph{Say it:} \emph{hādhihi masʾala thānawiyya.} — This is a secondary matter.
  \item \emph{Recall:} say the Arabic for negative, then the Arabic for basic, then say \emph{thānawī} again
\end{itemize}

\begin{tcolorbox}[breakable,colback=teal!4,colframe=teal!35!black,title={Before you move on}]
What does \ar{ثانوي} mean? (\textquotedblleft{}Secondary\textquotedblright{}.) Say it once more.
\end{tcolorbox}
\section[\texorpdfstring{\ar{شامل} (shāmil)}{شامل (shāmil)}]{\ar{شامل} (shāmil) — comprehensive}

\label{lesson:AR-C419-shamil}



\begin{quote}
Before the new one: say the Arabic for basic, then the Arabic for secondary.
\end{quote}

\subsection*{You'll want to know: \ar{شامل}}
\textbf{\ar{شامل}} — \emph{shāmil} — \textquotedblleft{}comprehensive\textquotedblright{}.

\begin{quote}
\ar{قرأت} \ar{تقريرا} \ar{شاملا} \ar{عن} \ar{المدينة}.
\end{quote}

\begin{quote}
\emph{qaraʾtu taqrīran shāmilan ʿan al-madīna.}
\end{quote}

\begin{quote}
I read a comprehensive report about the city.
\end{quote}

\begin{cousinweb}
\textbf{\ar{شامل}} (\emph{shāmil}) is built on the native Arabic root \emph{sh-m-l}, \textquotedblleft{}to comprise, to cover\textquotedblright{} (Lane). \emph{Shāmil} is the active participle of \emph{shamila}, to include: covering everything, comprehensive.
\end{cousinweb}

\begin{grammarlens}[title={What you've built}]
One more word for saying what happens if.
\end{grammarlens}

\subsection*{Your turn}
\begin{itemize}
\raggedright
  \item \emph{Say it:} \emph{shāmil}
  \item \emph{Say it:} \emph{shāmil}, once more
  \item \emph{Say it:} \emph{qaraʾtu taqrīran shāmilan ʿan al-madīna.} — I read a comprehensive report about the city.
  \item \emph{Recall:} say the Arabic for basic, then the Arabic for secondary, then say \emph{shāmil} again
\end{itemize}

\begin{tcolorbox}[breakable,colback=teal!4,colframe=teal!35!black,title={Before you move on}]
What does \ar{شامل} mean? (\textquotedblleft{}Comprehensive\textquotedblright{}.) Say it once more.
\end{tcolorbox}
